% ============================================
% Johns Hopkins Letter Template
% Uses the built-in 'letter' class
% ============================================

\documentclass[11pt,letterpaper]{letter}

\usepackage{graphicx}
\usepackage{eso-pic}
\usepackage{geometry}
\usepackage{microtype}
\usepackage[utf8]{inputenc}
\usepackage[hidelinks]{hyperref}

% ----- Page Setup -----
\geometry{left=25mm,right=25mm,top=25mm,bottom=25mm}

% ----- Logo Placement (foreground, first page only) -----
\newcommand{\PutJHULogoFG}[1]{%
  \AddToShipoutPictureFG*{%
    \AtPageUpperLeft{%
      \hspace{10mm}%
      \raisebox{-38mm}{%
        \includegraphics[width=6cm]{#1}%
      }%
    }%
  }%
}

% ----- Sender Information -----
\signature{Decory Edwards}
\address{Department of Economics \\ Johns Hopkins University \\ Baltimore, MD 21218 \\ \texttt{dedwar65@jh.edu}}

% ----- Recipient Info -----
\begin{document}
\PutJHULogoFG{Images/jhulogo.png} % show logo only on first page

\begin{letter}{Hiring Committee \\
Analysis Group \\
111 Huntington Avenue, 10th Floor \\
Boston, MA 02199}

\opening{Dear Hiring Committee,}

I am writing to express my interest in the Associate (PhD, 2026 Start Date) position at Analysis Group. I am a Ph.D.\ candidate in the Department of Economics at Johns Hopkins University (advisors: Christopher D.\ Carroll, Nicholas W.\ Papageorge, and Stelios Fourakis), and I expect to receive my degree in May 2026. My primary areas of specialization are macroeconomic modeling, wealth and income dynamics, and computational economics.

My research agenda focuses on the ability of macroeconomic models to generate distributions of wealth similar to those observed in the data. A key contributor to wealth accumulation is the rate of return earned on assets, so understanding the data on wealth holdings and how to compute reliable measures of returns is central to my work. These topics are closely aligned with Analysis Group’s focus on combining rigorous quantitative methods with real-world applications in litigation, policy, and business consulting.

In my job market paper, I match wealth moments from the Survey of Consumer Finances by simulating a model in which households differ in the return earned on their assets. This modeling choice is motivated by the growing empirical evidence documenting substantial heterogeneity in portfolio returns across individuals, even within narrowly defined asset classes. In another project, I use the Health and Retirement Study to analyze how trust relates to both income and portfolio returns. I replicate the well-known hump-shaped relationship between trust and income found in the literature, and show that a similar relationship holds for returns to wealth. I also characterize the distribution of the persistent component of returns to net worth across survey waves—findings that directly inform my structural modeling choices.

I am particularly drawn to Analysis Group’s interdisciplinary approach, where economists, data scientists, and subject-matter experts collaborate to tackle complex business and policy challenges. My background in building and calibrating quantitative models, managing large microdata sets, and producing transparent, reproducible analyses positions me well to contribute to empirical projects involving damages estimation, financial valuation, competition analysis, and policy evaluation.

I would be excited to join Analysis Group’s collaborative environment, where rigorous analysis is paired with mentorship and professional development. I believe my quantitative and computational skills would allow me to deliver clear, data-driven insights to clients and experts across practice areas.

Thank you for considering my application. I would welcome the opportunity to discuss how my research and analytical training could contribute to the impactful work done at Analysis Group.

\closing{Sincerely,}

\end{letter}
\end{document}
